\documentclass[11pt]{article}
\usepackage[utf8]{inputenc}
\usepackage[T1]{fontenc}
\usepackage[margin=1in]{geometry}
\usepackage{amsmath, amssymb}
\usepackage{booktabs}
\usepackage{array}
\usepackage{hyperref}

\begin{document}

\begin{center}
{\LARGE First Exam — Lessons 0 \& 1 (Version A)}\\[0.75em]
\end{center}

\noindent\textbf{Time allowed:} 60 minutes\\
\textbf{Resources:} Personal notes only. Electronic devices, calculators, and statistical tables are not permitted.\\
\textbf{Instructions:} Answer all four questions. Show algebraic steps and clearly label each part. Unless stated otherwise, provide exact values (fractions, radicals, exponentials). Partial credit is available when reasoning is clearly explained.

\vspace{0.75em}
\noindent\begin{tabular}{@{}p{2.2cm}p{8.2cm}r@{}}
\toprule
\textbf{Question} & \textbf{Topic} & \textbf{Points}\\
\midrule
1 & Probability structures and independence & 25\\
2 & Bayes' rule and conditional reasoning & 25\\
3 & Exponential model derivations & 25\\
4 & Interpreting summaries and convergence & 25\\
\midrule
\textbf{Total} & & \textbf{100}\\
\bottomrule
\end{tabular}

\vspace{1em}
\hrule
\vspace{1em}

\section*{Question 1 — Probability Foundations (25 points)}

A coin is tossed twice. Throughout parts (a)–(c), assume the coin is fair.

\begin{enumerate}
  \item (6 pts) Specify a probability space $\,(\Omega, \mathcal{F}, P)\,$ for this experiment. Clearly list the sample space, a natural $\sigma$-algebra, and the probability measure.
  \item (6 pts) Let $A$ be the event ``at least one head appears'' and $B$ the event ``exactly one head appears.'' Compute $P(A)$, $P(B)$, and $P(A \cap B)$.
  \item (5 pts) Are $A$ and $B$ independent? Justify your answer using the formal definition.
  \item (8 pts) Now suppose the coin has probability $p$ of landing heads, where $0 < p < 1$. With the same definitions of $A$ and $B$, derive expressions for $P_p(A)$, $P_p(B)$, and $P_p(A \cap B)$. For which value(s) of $p$ are $A$ and $B$ independent? Provide algebraic reasoning.
\end{enumerate}

\vspace{0.75em}
\hrule
\vspace{0.75em}

\section*{Question 2 — Bayes' Rule in Component Testing (25 points)}

A factory purchases circuit boards from two suppliers. Boards from $S_A$ account for $\tfrac{3}{5}$ of the total, and boards from $S_B$ account for the remaining $\tfrac{2}{5}$. Historical quality data show that $\tfrac{1}{10}$ of boards from $S_A$ are defective, while $\tfrac{1}{20}$ of boards from $S_B$ are defective. Each board is screened by a test that returns ``positive'' for a defective board with probability $\tfrac{9}{10}$, and returns ``positive'' for a non-defective board with probability $\tfrac{1}{20}$.

\medskip
\noindent\textit{Recall:} sensitivity $= P(\text{Positive} \mid \text{Defective})$ and specificity $= P(\text{Negative} \mid \text{Non-defective})$. The false positive rate is $1-\text{specificity} = P(\text{Positive} \mid \text{Non-defective})$.

\begin{enumerate}
  \item (7 pts) Compute the probability that a randomly selected board produces a positive test result.
  \item (8 pts) Given a positive test result, determine the probability that the board is defective. Express your answer as a reduced fraction.
  \item (6 pts) Given a negative test result, determine the probability that the board is defective. Express your answer as a reduced fraction.
  \item (4 pts) Let $\pi = P(\text{Defective})$, $s = P(\text{Positive} \mid \text{Defective})$, and $f = P(\text{Positive} \mid \text{Non-defective})$. Derive a formula for $P(\text{Defective} \mid \text{Positive})$ in terms of $\pi$, $s$, and $f$, and briefly state how decreasing $f$ affects this probability.
\end{enumerate}

\vspace{0.75em}
\hrule
\vspace{0.75em}

\section*{Question 3 — Exponential Waiting-Time Model (25 points)}

Let $T$ be the waiting time (in hours) until the next service request arrives, and suppose $T$ follows an exponential distribution with rate $\lambda = \tfrac{1}{4}$.

\medskip
\noindent\textbf{Recall:} The probability density function of an exponential random variable with rate $\lambda$ is given by
\[
f_T(t) =
\begin{cases}
\lambda e^{-\lambda t}, & t \ge 0,\\
0, & t < 0.
\end{cases}
\]

\begin{enumerate}
  \item (5 pts) Derive the cumulative distribution function $F_T(t)$.
  \item (5 pts) Using your result from part (1), compute $P(1 \le T \le 4)$ and leave the answer in exponential form.
  \item (9 pts) Compute $E[T]$ and $\operatorname{Var}(T)$ directly from integrals. Show the essential integration steps leading to exact values.
  \item (6 pts) Define $Y = 3T + 2$. Compute $E[Y]$ and $\operatorname{Var}(Y)$ using properties of expectation and variance. Provide exact values.
\end{enumerate}

\vspace{0.75em}
\hrule
\vspace{0.75em}

\section*{Question 4 — Interpreting Summaries and Convergence (25 points)}

Table~\ref{tab:a-sim} summarises the results of a simulation: for each sample size $n$, 1{,}000 independent samples of size $n$ were drawn from an exponential distribution with mean $4$, and the sample mean was recorded for each run.

\begin{table}[h]
\centering
\caption{Behaviour of Sample Means}\label{tab:a-sim}
\begin{tabular}{@{}rcc@{}}
\toprule
\textbf{Sample size $n$} & \textbf{Average of the 1,000 sample means} & \textbf{SD of the 1,000 sample means}\\
\midrule
5  & 4.08 & 1.80 \\
20 & 4.02 & 0.90 \\
80 & 4.01 & 0.44 \\
\bottomrule
\end{tabular}
\end{table}

Another analysis compared a standardised dataset (mean 0, variance 1) to the standard normal distribution. Selected points from the QQ-plot are listed below.

\begin{table}[h]
\centering
\caption{Selected QQ-Plot Quantiles}\label{tab:a-qq}
\begin{tabular}{@{}rc@{}}
\toprule
\textbf{Normal quantile} & \textbf{Sample quantile}\\
\midrule
-2.0 & -3.1 \\
-1.0 & -1.6 \\
 0.0 &  0.2 \\
 1.0 &  1.9 \\
 2.0 &  3.4 \\
\bottomrule
\end{tabular}
\end{table}

\medskip
Answer the following:
\begin{enumerate}
  \item (9 pts) Use Table~\ref{tab:a-sim} to explain how the Law of Large Numbers manifests in this simulation. Refer explicitly to the numerical values.
  \item (8 pts) Still relying on Table~\ref{tab:a-sim}, discuss how the variability of the sample mean changes with $n$ and connect your explanation to the Central Limit Theorem.
  \item (8 pts) Interpret Table~\ref{tab:a-qq}. Describe the shape of the corresponding QQ-plot, assess the plausibility of a normal model, and recommend a modelling or diagnostic action before relying on normal-based methods.
\end{enumerate}

\vspace{0.75em}
\hrule
\vspace{0.25em}

\noindent\textit{End of exam.} Ensure answers are clearly labelled and presented legibly.

\end{document}
